\section{The Fourth-Order Trotter Bound: Tighter Prefactors via BCH}
\label{sec:suzuki4}

The fourth-order Suzuki product formula $S_4(t)$ is the palindromic composition
of five Strang blocks with coefficients $(p, p, 1{-}4p, p, p)$, where
$p = 1/(4-4^{1/3})$ satisfies the Suzuki cubic $4p^3 + (1-4p)^3 = 0$:
\begin{equation}
  S_4(t)
    \;=\; S_2(pt)\,S_2(pt)\,S_2((1{-}4p)t)\,S_2(pt)\,S_2(pt),
  \qquad
  S_2(s) = e^{sA/2}\, e^{sB}\, e^{sA/2}.
\end{equation}
The Trotter error $\norm{S_4(t) - e^{t(A+B)}}$ vanishes to fourth order:
there exists a constant $C(A,B)$ such that
\begin{equation}
  \norm{S_4(t) - e^{t(A+B)}} \;\le\; C(A,B)\cdot t^5
\end{equation}
for all sufficiently small $t \ge 0$. The important question---answered
variously in the
literature with different prefactors---is: \emph{what is $C(A,B)$}?

Childs et al.~\cite{childsTheoryTrotterError2021} give, in the unitary
setting (skew-adjoint $A, B$, equivalently Hermitian generators evolving as
$e^{-itH}$) and for $t \ge 0$, an explicit bound of the
form
\begin{equation}\label{eq:childs-bound}
  \norm{S_4(t) - e^{t(A+B)}} \;\le\; t^5 \cdot \sum_{i=1}^{8} \alpha_i\, \norm{C_i(A,B)},
\end{equation}
where the $C_i$ are the eight four-fold nested commutators
$[X_1,[X_2,[X_3,[B,A]]]]$ with $X_j \in \{A, B\}$, and the coefficients $\alpha_i$
range from $0.0046$ to $0.0284$. This section discusses what we can (and cannot)
say about the tightness of these prefactors, how a Computer-Algebra-System
(CAS)-assisted BCH expansion yields strictly smaller coefficients, and how the
result is integrated with our Lean formalization.

\subsection{Childs et al.'s prefactors: rigorous, but not proven tight}
\label{sec:childs-heuristic}

A point of precision first: \cref{eq:childs-bound} is a \emph{rigorous}
upper bound. It is stated and proved as a proposition (Prop.~J.1 in the
arXiv version of~\cite{childsTheoryTrotterError2021}), derived via a
Duhamel expansion and the triangle inequality, and it holds for all
$t \ge 0$ in the unitary setting. What Childs et al.\ describe as
``heuristic'' are the \emph{strategies} used to make the prefactors small,
and what they leave unproven is tightness: ``we do not have a rigorous
proof of the tightness of these bounds.''

Concretely, they derive \cref{eq:childs-bound} by an elegant but
``balanced factoring,'' whose optimality is not proved, of the twelve-factor Duhamel
representation of the $S_4$ error integrand. The 4-fold commutator structure
emerges naturally from this Duhamel form; the specific numerical values of
$\alpha_i$ come from bounding the factored integrals term by term.

Two independent sources of looseness enter this derivation:

\paragraph{Over-completeness of the $C_i$ basis.}
The weight-five free Lie algebra on two generators has dimension six (Witt's
formula: $6 = \tfrac{1}{5}\sum_{d \mid 5} \mu(5/d)\,2^d$). Childs's eight
commutators $C_1,\ldots,C_8$ are therefore an \emph{over-complete} spanning
set: eight vectors in a six-dimensional space, satisfying two independent
linear relations. Whenever a Lie-algebra element such as the exact residual
$R_5$ is written as $\sum_i \beta_i C_i$, the signed
coefficients $\beta_i$ are not unique: one has two degrees of freedom
(corresponding to the two relations) to redistribute mass across the $C_i$'s.
Childs et al.'s $\alpha_i$ are \emph{not} coordinates of $R_5$ in this (or
any) representation---they bound the terms of the factored Duhamel integral
individually---so over-completeness enters their derivation only implicitly,
through which nested commutators the factoring happens to produce. For a
BCH-style bound it enters directly: each admissible representation yields a
bound $\sum_i \abs{\beta_i}\norm{C_i}$, and minimizing this weighted,
instance-dependent objective over representations (not an ordinary
coefficient $\ell^1$ norm) would generally give something smaller than any
fixed choice, including ours.

\paragraph{Balanced factoring vs.\ exact expansion.}
Childs's derivation bounds the Duhamel integral \emph{without} computing the
explicit Baker--Campbell--Hausdorff (BCH) expansion of $\log S_4(t)$. Each
term of the balanced factoring is bounded individually via the triangle
inequality. An exact BCH expansion gives a single expression for $R_5$ with
\emph{signed} coefficients whose cancellations have already been applied
(leveraging the Suzuki cubic condition and palindromic symmetry). Bounding
an expression with cancellations already applied is tighter than bounding
individual terms before cancellation.

Neither source of looseness affects validity---the $\alpha_i$ bound is
correct as stated---but both leave room below it, and that room is exactly
what the BCH expansion of \cref{sec:bch-expansion} occupies.

\subsection{The BCH expansion and its cancellations}
\label{sec:bch-expansion}

The BCH expansion of $\log S_4(t) - t(A+B)$ has only odd-order terms (by
palindromic symmetry), of which the leading one cancels under the Suzuki
condition:
\begin{equation}\label{eq:log-s4}
  \log S_4(t) - t(A+B) \;=\; \underbrace{0}_{O(t^2)} + \underbrace{0}_{O(t^3)}
  + \underbrace{0}_{O(t^4)} + t^5\,R_5(A,B) + O(t^7).
\end{equation}

\paragraph{$O(t^2)$ and $O(t^4)$ vanish by palindromic symmetry.}
The palindromic symmetry of $S_4$---invariance under time reversal
$\bar S_4(t) := S_4(-t)^{-1} = S_4(t)$---forces all even-order BCH
coefficients to vanish identically, regardless of $p$. This is Yoshida's
theorem for symmetric compositions.

\paragraph{$O(t^3)$ vanishes under the Suzuki cubic.}
The third-order BCH coefficient of $S_4$ has the form
$\sum_i c_i^3\, E_3(A,B)$, where $c_i \in \{p, p, 1{-}4p, p, p\}$ are the
block coefficients and
\begin{equation}\label{eq:E3}
  E_3 \;=\; \tfrac{1}{12}\comm{B}{\comm{B}{A}} \;-\; \tfrac{1}{24}\comm{A}{\comm{A}{B}}
\end{equation}
is the Strang BCH cubic coefficient. The Suzuki parameter $p = 1/(4-4^{1/3})$ is
precisely the root of
\begin{equation}
  \sum_{i=1}^{5} c_i^3 \;=\; 4p^3 + (1-4p)^3 \;=\; 0.
\end{equation}
Hence the $O(t^3)$ coefficient vanishes.

\paragraph{The $O(t^5)$ residual.}
The fifth-order residual $R_5(A,B)$ lives in the six-dimensional weight-5
Lie algebra. It is a \emph{specific} Lie polynomial, computable in closed
form from the BCH expansion of five palindromic exponentials.

\subsection{CAS computation and Lean integration}
\label{sec:cas}

Closed-form computation of $R_5$ requires tracking $\sim 30$ five-letter
monomials in $A$, $B$ through the BCH series---feasible but tedious by hand,
and natural for a Computer Algebra System.

We use a \textbf{symbolic non-commutative free algebra over $\mathbb{Q}[p]$},
implemented in Python with \texttt{sympy} (see
\leanref{../scripts/compute\_bch\_prefactors.py}{1}{scripts/compute\_bch\_prefactors.py}
in the repository):
\begin{enumerate}
  \item Elements of the free algebra are dictionaries mapping words
    (tuples over $\{A, B\}$) to rational-function coefficients in
    $\mathbb{Q}[p]$.
  \item $S_2(s)$ is computed by truncated power series for $e^{(s/2)A}$,
    $e^{sB}$, $e^{(s/2)A}$, composed non-commutatively.
  \item $S_4(t)$ is computed as the product of five $S_2$ blocks with
    coefficients $(p, p, 1{-}4p, p, p)$, truncated to total degree $5$.
  \item $\log S_4(t) - t(A+B)$ is computed by the truncated series
    $\log(1 + y) = y - y^2/2 + y^3/3 - y^4/4 + y^5/5$ with
    $y = S_4(t) - 1$.
  \item The Suzuki cubic reduction $p^3 = (48p^2 - 12p + 1)/60$
    is applied to every coefficient, reducing all powers of $p$ to
    polynomials of degree~$\le 2$.
  \item The script verifies degrees $2$, $3$, $4$ vanish identically and
    extracts the degree-5 part.
  \item The degree-5 residual is projected onto the Childs
    basis $\{C_1, \ldots, C_8\}$ via Gauss--Jordan elimination.
    Since the basis is over-complete, the projection has two free
    parameters; the script chooses the projection setting both to zero.
  \item Specializing at $p = 1/(4-4^{1/3}) \approx 0.4145$ gives numerical
    values for the eight signed coefficients $\beta_i(p)$; the Lean
    development stores the certified rational ceilings
    $\gamma_i \ge \abs{\beta_i(p)}$ at the $10^{-6}$ grid.
\end{enumerate}

\paragraph{The symbolic result.}
Before numerical specialization, the eight signed coefficients $\beta_i(p)$
are polynomials of degree~$\le 2$ in $p$:
\begin{equation}\label{eq:bch-gamma}
  \setlength{\arraycolsep}{1pt}
  \begin{array}{rlrlrl}
    \beta_1(p) &= \tfrac{127}{144000}\,p^2 + \tfrac{13}{36000}\,p - \tfrac{1}{24000}, \qquad &
    \beta_5(p) &= \tfrac{31}{9000}\,p^2 - \tfrac{13}{18000}\,p + \tfrac{1}{12000}, \\[3pt]
    \beta_2(p) &= \tfrac{1}{12000}\,p^2 + \tfrac{13}{6000}\,p - \tfrac{1}{4000}, &
    \beta_6(p) &= \tfrac{31}{3000}\,p^2 - \tfrac{13}{6000}\,p + \tfrac{1}{4000}, \\[3pt]
    \beta_3(p) &= 0, &
    \beta_7(p) &= 0, \\[3pt]
    \beta_4(p) &= -\tfrac{61}{9000}\,p^2 + \tfrac{13}{3000}\,p - \tfrac{1}{2000}, &
    \beta_8(p) &= \tfrac{1}{18000}\,p^2 + \tfrac{13}{9000}\,p - \tfrac{1}{6000}.
  \end{array}
\end{equation}
Note that $\beta_3$ and $\beta_7$ are identically zero---a consequence of
the particular projection onto the over-complete basis. Throughout,
$\beta_i(p)$ denotes the exact signed coefficients and $\gamma_i$ the
certified nonnegative rational ceilings $\gamma_i \ge \abs{\beta_i(p)}$
stored in the Lean development.

\paragraph{Comparison with Childs.}
Specializing at Suzuki $p$ and rounding up to the certified ceilings
$\gamma_i$, all non-zero $\gamma_i$ are strictly
smaller than the corresponding Childs coefficient $\alpha_i$, most by
an order of magnitude or more (\cref{tab:bch-vs-childs}).

\begin{table}[ht]
\centering
\begin{tabular}{@{}clrrl@{}}
\toprule
\textbf{Index} & \textbf{Commutator $C_i$} & \textbf{Ceiling~$\gamma_i$} & \textbf{Childs $\alpha_i$} & \textbf{Ratio} \\
\midrule
$C_1$ & $[A,[A,[A,[B,A]]]]$ & $0.000260$ & $0.0047$ & $18\times$ \\
$C_2$ & $[A,[A,[B,[B,A]]]]$ & $0.000663$ & $0.0057$ & $8.6\times$ \\
$C_3$ & $[A,[B,[A,[B,A]]]]$ & $0$        & $0.0046$ & $\infty$ \\
$C_4$ & $[A,[B,[B,[B,A]]]]$ & $0.000132$ & $0.0074$ & $56\times$ \\
$C_5$ & $[B,[A,[A,[B,A]]]]$ & $0.000376$ & $0.0097$ & $26\times$ \\
$C_6$ & $[B,[A,[B,[B,A]]]]$ & $0.001128$ & $0.0097$ & $8.6\times$ \\
$C_7$ & $[B,[B,[A,[B,A]]]]$ & $0$        & $0.0173$ & $\infty$ \\
$C_8$ & $[B,[B,[B,[B,A]]]]$ & $0.000442$ & $0.0284$ & $64\times$ \\
\bottomrule
\end{tabular}
\caption{Certified rational ceilings $\gamma_i \ge \abs{\beta_i(p)}$ vs.\ Childs et al.'s
coefficients $\alpha_i$ (at Suzuki $p = 1/(4-4^{1/3})$). The ``Ratio''
column is $\alpha_i/\gamma_i$---how many times tighter the certified
coefficient is. Entries $C_3$ and $C_7$ vanish exactly in our chosen
projection; other valid projections would redistribute their contribution
across the remaining six.}
\label{tab:bch-vs-childs}
\end{table}

\paragraph{Lean integration.}
The script emits a Lean snippet for the
\texttt{bchTightPrefactors} definition in
\leansrc{Suzuki4ViaBCH.lean}. The corresponding Lean theorem
\texttt{norm\_suzuki4\_level3\_explicit} states that there exist
$\delta > 0$ and $K' \ge 0$ such that, for all $0 \le t < \delta$,
\begin{equation}\label{eq:level3-sharp}
  \norm{S_4(t) - e^{t(A+B)}}
    \;\le\; t^5 \cdot \sum_{i=1}^{8} \gamma_i\, \norm{C_i}
    \;+\; K' \, t^6 .
\end{equation}
The leading coefficient is the CAS-certified $\gamma$-sum \emph{as it appears in
the statement}; only the $O(t^6)$ remainder is existentially quantified. The
theorem is fully proved at the project level: \texttt{\#print axioms
norm\_suzuki4\_level3\_explicit} returns only \texttt{[propext,
Classical.choice, Quot.sound]}, and it requires no C*-algebra or anti-Hermitian
structure---any complete normed algebra suffices.

The intermediate bridge theorem
\texttt{bch\_w4Deriv\_level3\_tight}---historically axiomatized as the
BCH-derived pointwise residual---is now a theorem composing the
five-factor palindromic quintic identification from the companion
\href{https://github.com/Jue-Xu/Lean-BCH}{Lean-BCH} project
(\texttt{suzuki5\_log\_product\_quintic\_tight\_at\_suzukiP}, revision
\texttt{05e8c52}) with the per-$i$ numerical bounds
$\abs{\beta_i(\text{suzukiP})} \le \gamma_i$ proved by \texttt{nlinarith}
on the tight rational interval $41449/100000 < \text{suzukiP} <
41450/100000$.

A companion theorem \texttt{norm\_suzuki4\_level3\_le\_childs\_pointwise}
formally verifies that the BCH-derived leading coefficient is dominated \emph{by}
Childs's---i.e.\ that our bound is the tighter of the two:
\begin{equation}
  t^5 \cdot \sum_{i} \gamma_i\,\norm{C_i}
    \;\le\;
  t^5 \cdot \sum_{i} \alpha_i\,\norm{C_i}.
\end{equation}
The element-wise comparison $\gamma_i \le \alpha_i$ is verified by
\texttt{nlinarith} from the explicit numerical values. Substituting it into
\cref{eq:level3-sharp} gives \texttt{norm\_suzuki4\_childs\_explicit}:
Childs et al.'s coefficients in the statement, still carrying the same
$K' t^6$ remainder---the Level~1 form described below. Sharpening further,
\texttt{norm\_suzuki4\_le\_childs\_near\_zero} shows that whenever the two
leading coefficients differ strictly, the BCH bound beats Childs et al.'s
\emph{without any remainder term} on a neighbourhood of the origin (the
radius $\delta$ is assembled from the existential Level-3 witnesses
$\delta_0, K'$, hence not explicit):
$\norm{S_4(t) - e^{t(A+B)}} \le t^5 \sum_i \alpha_i \norm{C_i}$ for
$0 \le t < \delta$.

\subsection{A local small-\texorpdfstring{$\tau$}{tau} refinement at order \texorpdfstring{$\tau^7$}{tau7}: the \texorpdfstring{$R_7$}{R7} correction}
\label{sec:r7}

A subtle question: is our leading-order Level 3 bound $t^5 \sum \gamma_i \norm{C_i}$
actually tighter than Childs's uniform bound, or does the comparison depend
on how we handle higher-order BCH contributions?

\paragraph{What the comparison does and does not establish.}
A natural worry is that Childs et al.'s larger coefficients $\alpha_i$ might
be \emph{uniformly tight} on some finite interval $[0, t^*]$, so that
$\alpha_i \cdot \norm{C_i}$ is the correct envelope over that interval
(folding in higher BCH orders), while our leading $\gamma_i$ is only an
asymptotic $t \to 0$ value. Two observations bear on this---neither of
which settles it:

\begin{enumerate}[leftmargin=*]
  \item \emph{The leading-order gap is large.} The ratio
    $\alpha_i / \gamma_i$ ranges from $8.6\times$ to $64\times$
    (\cref{tab:bch-vs-childs}). For the $\alpha_i$ to be uniformly tight on
    $[0, t^*]$, higher-order BCH contributions at $t = t^*$ would need
    to exactly saturate the gap $\alpha_i - \gamma_i$, and
    no mechanism in the balanced-factoring derivation---a
    triangle-inequality
    envelope of 12 individual Duhamel terms---targets this kind of
    saturation.

  \item \emph{Tightness is unproven, not disproven.} The original paper's
    statement ``we do not have a rigorous proof of the tightness of these
    bounds'' leaves uniform tightness \emph{open}, in both directions.
\end{enumerate}

What our development proves is confined to the leading coefficient
(termwise $\gamma_i \le \alpha_i$, so the $\gamma$-sum never exceeds the
$\alpha$-sum) and, under the strict-gap hypothesis, to a neighbourhood of
$0$ on which the full BCH bound undercuts the $\alpha$-envelope
(\texttt{norm\_suzuki4\_le\_childs\_near\_zero}). Global or uniform
optimality---of either family of coefficients---remains unknown.

\paragraph{A uniform BCH bound via \texorpdfstring{$R_7$}{R7}.}
To produce a rigorously uniform bound (not merely asymptotic), we extend
the CAS computation to order~7 (see
\leanref{../scripts/compute\_bch\_r7.py}{1}{compute\_bch\_r7.py}). The
script confirms:
\begin{itemize}[nosep]
  \item Orders $2$, $3$, $4$, and $6$ all vanish identically under
    palindromic symmetry + Suzuki cubic (extending the cancellations
    verified at order 5).
  \item Order $7$ has a non-zero residual $R_7(A, B)$ which, at Suzuki $p$,
    expands to 126 non-zero 7-letter words.
\end{itemize}
A straightforward triangle-inequality bound over these 126 words, using
$\norm{w} \le \max(\norm{A}, \norm{B})^7$ for any 7-letter word $w$, yields
the explicit uniform coefficient
\begin{equation}\label{eq:r7-K}
  K := \sum_{w} \abs{\mathrm{coef}(w)} \;\approx\; 0.01951,
  \qquad
  \norm{R_7(A, B)} \le K \cdot \max(\norm{A}, \norm{B})^7.
\end{equation}
The per-AB-degree breakdown of $K$ is shown in \cref{tab:r7-breakdown}.

\begin{table}[ht]
\centering
\begin{tabular}{@{}lrlr@{}}
\toprule
\textbf{AB signature} & $K_k$ \\
\midrule
$\norm{A}^1 \norm{B}^6$ & $0.000828$ \\
$\norm{A}^2 \norm{B}^5$ & $0.003835$ \\
$\norm{A}^3 \norm{B}^4$ & $0.006134$ \\
$\norm{A}^4 \norm{B}^3$ & $0.005510$ \\
$\norm{A}^5 \norm{B}^2$ & $0.002662$ \\
$\norm{A}^6 \norm{B}^1$ & $0.000540$ \\
\midrule
\textbf{Total $K$} & $0.019509$ \\
\bottomrule
\end{tabular}
\caption{R$_7$ uniform bound $K = \sum_k K_k$ split by AB signature. The
symmetric profile ($K_3, K_4$ peak) reflects palindromic structure.}
\label{tab:r7-breakdown}
\end{table}

\paragraph{The Level 4 existential-$\delta$ bound.}
Combining the certified ceilings $\gamma_i$ of the $R_5$ coefficients with
the $R_7$ uniform constant $K$, we obtain a bound on a small interval
$[0, \delta)$. The interval is a limitation of the proof route, not of the
statement: the bound is extracted from the BCH identification of
$\log S_4(t)$, which is available only within the convergence radius of the
BCH series, and $\delta$ records that radius. Duhamel-type analyses reach
all $t \ge 0$ directly---Childs et al.'s Prop.~J.1 is global in $t$ in the
unitary setting, and our compounded total-error bounds
(\cref{sec:s4-convergence}) hold, for every fixed $t$, for all sufficiently
large $n$ in a general normed algebra, at the price of an exponential
growth factor.
\begin{equation}\label{eq:level4}
\begin{aligned}
  &\exists\, \delta > 0,\ \exists\, C \ge 0,\ \forall\, t \in [0, \delta), \\
  &\quad \norm{S_4(t) - e^{t(A+B)}}
    \;\le\;
  C \cdot \Big( t^5 \cdot \sum_{i=1}^{8} \gamma_i \norm{C_i}
  \;+\;
  t^7 \cdot K \cdot \max(\norm{A}, \norm{B})^7
  \;+\;
  t^8 \Big).
\end{aligned}
\end{equation}
The multiplier $C$ absorbs the exp-Lipschitz inflation
$\exp(\norm{tH} + \norm{\delta_{\mathrm{bch}}})$ from lifting the BCH
log identification to $S_4$; the $t^8$ tail absorbs the higher-order
remainder of the BCH expansion ($R_9, R_{11}, \ldots$). For the small-$t$
regime of Trotter splitting these are negligible; if desired the same
pipeline can extend the CAS to arbitrary order with correspondingly
smaller corrections.

The Lean formalization adds the uniform constant
\leanref{Suzuki4ViaBCH.lean}{1116}{bchR7UniformConstant}
and the bound
\begin{equation}
  \norm{R_7} \;\le\; K \cdot \max(\norm{A}, \norm{B})^7,
\end{equation}
encoded as \leanref{Suzuki4ViaBCH.lean}{1122}{bchR7Bound}.
The Lean-BCH bridge corollary
\texttt{suzuki5\_log\_product\_septic\_at\_suzukiP} (now a theorem
upstream, with its $\tau^7$ Childs-7-basis identification factored
through two deeper Lean-BCH septic theorems at the pinned revision;
see \cref{sec:caveats}) composes with exp-Lipschitz on the Lean-Trotter
side to yield the theorem
\begin{leancode}
theorem norm_suzuki4_level4_uniform (A B : 𝔸)
    (hA : star A = -A) (hB : star B = -B) :
    let p : ℝ := 1 / (4 - (4 : ℝ) ^ ((1 : ℝ) / 3))
    ∃ δ > 0, ∃ C ≥ 0, ∀ τ : ℝ, 0 ≤ τ → τ < δ →
      ‖suzuki4Exp A B p τ - exp (τ • (A+B))‖ ≤
        C * (τ^5 * bchTightPrefactors.boundSum A B +
             τ^7 * bchR7Bound A B +
             τ^8)
\end{leancode}

\paragraph{Where the Childs comparison actually lives.}
It is tempting to compare \cref{eq:level4} against Childs's
$t^5 \sum \alpha_i \norm{C_i}$ directly, but \cref{eq:level4} is the
wrong place to do it: its multiplier $C$ is existentially quantified,
so the quantity $C \sum_i \gamma_i \norm{C_i}$ that sits in front of
$t^5$ is not a number the statement pins down. The comparison belongs
at Level~3, where \cref{eq:level3-sharp} certifies the leading
coefficient outright and only the $t^6$ remainder is existential.
Absorbing that remainder into the $\alpha$--$\gamma$ gap yields Childs's
\emph{own} bound with \emph{no} remainder term on a neighbourhood of the
origin, which is exactly
\leanref{Suzuki4ViaBCH.lean}{1387}{norm\_suzuki4\_le\_childs\_near\_zero}
(stated in \cref{sec:cas}): write $K' t^6 = t^5 \cdot (K' t)$ and observe
that $K' t$ falls below the gap $\sum_i (\alpha_i - \gamma_i) \norm{C_i}$
as soon as $t < \mathrm{gap}/(K'+1)$.

Its hypothesis, $\sum_i \gamma_i \norm{C_i} < \sum_i \alpha_i \norm{C_i}$,
compares two \emph{concrete} reals computed from the eight commutator norms,
so it is checkable for any given $A, B$; since $\gamma_i \le \alpha_i$
termwise, it fails only in the degenerate case where every commutator with
$\gamma_i < \alpha_i$ vanishes---e.g.\ $\comm{A}{B} = 0$, where both bounds
are $0$ and there is nothing to beat. And because it routes through Level~3
rather than Level~4, it needs neither anti-Hermiticity nor a C*-algebra.

\paragraph{What Level 4 is, and is not.}
L4 is not the bound one plugs a number into. Both $C$ and $\delta$ are
existentially quantified in its statement, so at a given step size $t$ one
knows neither whether $t < \delta$ nor what $C$ is; and since $C$ multiplies
the leading term as well, L4 certifies \emph{less} about the $t^5$ coefficient
than \cref{eq:level3-sharp} does. For comparing certified leading
coefficients---where the
$8.6$--$64\times$ reduction of the $\gamma_i$ error prefactors \emph{would}
translate into a fourth-root reduction in the required fourth-order steps
in the small-$t$ regime---the
statement to use is the sharp Level~3 bound \cref{eq:level3-sharp}, whose
leading coefficient is the CAS-certified $\gamma$-sum verbatim. Since
$\delta$ and $K'$ there (and $N$, $K'$ in the compounded bounds of
\cref{sec:s4-convergence}) remain existential, what is certified is an
explicit asymptotic leading coefficient, not an end-to-end numerical step or
gate count; extracting the latter would require re-proving those statements
with explicit constants.

What L4 does contribute is evidence that the CAS-assisted BCH pipeline
\emph{iterates}. It carries the order-$\tau^7$ residual through the full
chain---Lean-BCH septic identification, Childs-7-basis projection,
exp-Lipschitz lift---and lands the explicit constant
$K = 0.01951$ of \cref{eq:r7-K} inside a formal statement, with the next
correction pushed out to $\tau^8$. An analogous $R_9$ extension is mechanical.
The existential $\delta$ is a limitation of the logarithmic proof route
rather than of the statement: it records the convergence radius within which
the BCH identification of $\log S_4(t)$ is available, and it is not forced
by the conclusion---Childs et al.'s Duhamel-based Prop.~J.1 holds for all
$t \ge 0$ in the unitary setting, and the compounded bounds of
\cref{sec:s4-convergence} hold in a general normed algebra with an
exponential growth factor. (A polynomial right-hand side valid for all
$t \ge 0$ cannot be expected in a general normed algebra---the error may
itself grow exponentially in $t$, though it need not: it vanishes
identically for commuting $A, B$ and stays bounded in the unitary
setting---so global bounds there carry exponential factors.) The
existential $C$, by contrast, is an artifact
of how the exp-Lipschitz inflation was packaged, and is precisely what
\cref{eq:level3-sharp} avoids.

\subsection{Summary of the four BCH-based bounds}
\label{sec:summary-levels}

\cref{tab:levels} summarizes the four formalized bounds.

\begin{table}[ht]
\centering
\footnotesize
\setlength{\tabcolsep}{4pt}
\begin{tabular}{@{}l@{\hspace{4pt}}l@{\hspace{4pt}}l@{\hspace{4pt}}p{4.2cm}@{}}
\toprule
\textbf{Level} & \textbf{Lean theorem} & \textbf{Form} & \textbf{Notes} \\
\midrule
\textbf{L1 (Childs)}
  & \texttt{norm\_suzuki4\_le\_childs\_near\_zero}
  & $t^5 \sum \alpha_i \norm{C_i}$
  & Childs's own coefficients, with \emph{no} remainder term, on $[0,\delta)$.
    Corollary of L3 via $\gamma_i \le \alpha_i$; no axiomatization of
    the Childs bound.
    Needs the gap hypothesis
    $\sum \gamma_i\norm{C_i} < \sum \alpha_i\norm{C_i}$. \\
\midrule
\textbf{L2 (BCH unit)}
  & \texttt{norm\_suzuki4\_level2\_explicit}
  & $t^5 \sum \norm{C_i} + K' t^6$
  & Unit coefficients from primitive BCH
    $\abs{\beta_i(\text{Suzuki-}p)} \le 1$. Holds for every $p$
    satisfying the cubic condition. \\
\midrule
\textbf{L3 (BCH leading)}
  & \texttt{norm\_suzuki4\_level3\_explicit}
  & $t^5 \sum \gamma_i \norm{C_i} + K' t^6$
  & Certified $\gamma_i$, smaller than Childs's coefficient-by-coefficient
    (8.6--64$\times$); leading-order only. \\
\midrule
\textbf{L4 (BCH septic)}
  & \texttt{norm\_suzuki4\_level4\_uniform}
  & $C\,(t^5\!\sum \gamma_i \norm{C_i} + t^7 K M^7 + t^8)$
  & Small-$\tau$ bound, $M = \max(\norm{A}, \norm{B})$. $C$ absorbs the
    exp-Lipschitz inflation; being existential, it makes the certified
    $t^5$ coefficient weaker than L3's. \\
\bottomrule
\end{tabular}
\caption{Four formalized BCH-based bounds for $\norm{S_4(t) - e^{t(A+B)}}$.
Each row names the theorem whose \emph{type}
carries the stated Form; all are $\exists \delta > 0$ statements on $[0,\delta)$,
with $K'$ (and, in L4, $C$) existentially quantified but bound \emph{before} $\tau$.
L1 is a \emph{corollary} of L3---no axiomatization of the Childs bound on
Trotter's side---while L2/L3/L4 are Lean theorems composing the corresponding
Lean-BCH bridge corollaries with exp-Lipschitz lifts. At Lean-BCH
revision \texttt{05e8c52} (2026-07-28), all four levels are axiom-free
(\texttt{\#print axioms} returns only the standard foundational axioms),
including L4's $\tau^7$ extension.
Hypotheses also differ: L1--L3 hold in any complete normed algebra with
$\norm{1} = 1$, whereas L4 requires anti-Hermitian $A, B$ in a
C\textsuperscript{*}-algebra (\cref{sec:s4-convergence}).
}
\label{tab:levels}
\end{table}

\subsection{From step error to convergence:
  \texorpdfstring{$O(1/n^4)$}{O(1/n4)}}
\label{sec:s4-convergence}

Levels~1--4 all bound the \emph{single-step} error
$\norm{S_4(t) - e^{t(A+B)}}$. What a simulation pays for, however, is the
\emph{total} error after $n$ steps of size $t/n$, and the two are separated by
the telescoping amplification of \cref{sec:bg-trotter}: a step error of order
$\tau^{k+1}$ compounds to a total error of order $1/n^{k}$. Formalizing that
passage completes the hierarchy---the first-order and Strang formulas already
carry total-error theorems (\cref{sec:assembly,sec:strang}), whereas $S_4$
carried only a step bound.

\paragraph{The bound.}
Let $p$ satisfy the Suzuki cubic condition $4p^3 + (1-4p)^3 = 0$, as the
standard $p = 1/(4 - 4^{1/3})$ does. Then for all $A, B$ in a complete normed
algebra and every evolution time $t$, there exist a constant $C > 0$ and a
threshold $N$ such that
\begin{equation}\label{eq:s4-quartic}
  \norm{\bigl(S_4(t/n)\bigr)^n - e^{t(A+B)}} \;\le\; \frac{C}{n^4}
  \qquad \text{for all } n \ge N,
\end{equation}
and consequently $\bigl(S_4(t/n)\bigr)^n \to e^{t(A+B)}$ as $n \to \infty$
(\leanref{Suzuki4Convergence.lean}{158}{suzuki4\_total\_error\_quartic},
\leanref{Suzuki4Convergence.lean}{257}{suzuki4\_convergence\_quartic}).
Both are axiom-free. At the standard $p$ the cubic hypothesis is discharged
internally, leaving a statement with no hypotheses beyond the ambient algebra
(\leanref{Suzuki4Convergence.lean}{326}{suzuki4\_convergence\_quartic\_suzukiP}).

\paragraph{Proof.}
Telescoping (\cref{sec:telescoping}) gives
$\norm{X^n - Y^n} \le n\,\norm{X-Y}\,\max(\norm{X},\norm{Y})^{n-1}$
with $X = S_4(t/n)$ and $Y = e^{(t/n)(A+B)}$, and two ingredients feed it.
The step error $\norm{X-Y} \le C_0\,\abs{t/n}^5$ is the single-step BCH bound
(\leanref{Suzuki4BchBound.lean}{386}{exists\_norm\_s4Func\_sub\_exp\_le\_t5}),
valid once $\abs{t/n} < \delta$---which is exactly what the threshold $N$
enforces. The damping factor is a \emph{constant}: writing $K$ for the sum of
the magnitudes of the eleven $S_4$ coefficients weighted by $\norm{A}$ and
$\norm{B}$, we prove $\norm{S_4(\tau)} \le e^{\abs{\tau} K}$
(\leanref{Suzuki4Convergence.lean}{106}{norm\_suzuki4Exp\_le}), whence
\[
  \max(\norm{X},\norm{Y})^{n-1} \;\le\; \bigl(e^{\abs{t/n}\,K}\bigr)^{n}
  \;=\; e^{\abs{t}\,K},
\]
independent of $n$. The exponential target composes exactly,
$\bigl(e^{(t/n)(A+B)}\bigr)^n = e^{t(A+B)}$, so no error enters from that side,
and the pieces combine to
$n \cdot C_0 \abs{t}^5 n^{-5} \cdot e^{\abs{t} K} = O(1/n^4)$.

\paragraph{Scope of the bound.}
In \cref{eq:s4-quartic} both $C$ and $N$ are existential: $N$ marks where the
step size enters the BCH regime $\abs{\tau} < \delta$, and $C$ is inherited
from the crude single-step bound, whose constant is expressed through
$\norm{A}$ and $\norm{B}$. That establishes the \emph{order} but not the
\emph{constant}. Two things distinguish it from the generic $O(1/n^2)$ rate for
$S_4$ (\leanref{Suzuki4.lean}{376}{suzuki4\_error\_rate\_sq}), which holds for
\emph{every} $n \ge 1$ and needs \emph{no} cubic condition: fourth order is
precisely what the hypothesis $4p^3 + (1-4p)^3 = 0$ buys, and it is bought only
beyond the threshold $N$. (Both constants are existentially quantified; the
$O(1/n^2)$ proof's witness, $19 s^3 e^{6s} + 1$ with $s = \norm{A} + \norm{B}$,
is visible in the source but not in the statement.)

\paragraph{Commutator scaling in the total error.}
Compounding the \emph{sharp} step bound \eqref{eq:level3-sharp} through the same
telescope---rather than the crude one---keeps the certified prefactors alive all
the way to the $n$-step statement. At the Suzuki parameter, for $t \ge 0$ there
are $K' \ge 0$ and $N$ such that, for all $n \ge N$,
\begin{equation}\label{eq:s4-quartic-comm}
  \norm{S_4(t/n)^n - e^{t(A+B)}}
    \;\le\;
  e^{tK} \cdot \Bigl(\sum_{i=1}^{8} \gamma_i \norm{C_i}\Bigr)\,
    \frac{t^5}{n^4}
  \;+\; \frac{K'}{n^5},
  \qquad K = \texttt{s4Rate} + \norm{A} + \norm{B},
\end{equation}
and the same holds with Childs's $\alpha_i$ in place of $\gamma_i$
(\leanref{Suzuki4TightConvergence.lean}{59}{suzuki4\_total\_error\_commutator\_scaling},
\texttt{suzuki4\_total\_error\_childs\_scaling}; both axiom-free, and again with
no star structure). This is the commutator-scaling statement in the sense of
Childs et al.: the coefficient of $1/n^4$ is a sum of \emph{nested-commutator
norms}, so it collapses as $A$ and $B$ approach commutativity---behaviour that
an $\exists C$ bound such as \cref{eq:s4-quartic} cannot express.

\paragraph{No star structure is needed.}
\Cref{eq:s4-quartic} and the Level~1--3 step bounds alike hold in \emph{any}
complete normed algebra with $\norm{1} = 1$: they reach $S_4$ through BCH and
an exp-Lipschitz estimate, and never invoke the anti-Hermitian isometry
$\norm{e^{sX}} = 1$. Their Lean statements assume only \texttt{NormedRing},
\texttt{NormedAlgebra} $\mathbb{R}$, \texttt{NormOneClass} and
\texttt{CompleteSpace}. Among the $S_4$ results, only the Level~4 uniform bound
genuinely requires the C\textsuperscript{*}/anti-Hermitian setting---it carries
the hypotheses $A^* = -A$ and $B^* = -B$ explicitly---as does the Duhamel track
of \cref{sec:commutator}, where the isometry is what makes the FTC argument
close. The tight 4th-order Trotter bound is therefore available for arbitrary
bounded generators, not only for the skew-adjoint generators of unitary
dynamics.

\paragraph{One operator, two spellings.}
$S_4$ enters the development twice: as five Strang steps with a built-in $1/n$
(\texttt{suzuki4Step}, the object carrying the generic $O(1/n^2)$ rate), and as
eleven exponentials in a free step size $\tau$ (\texttt{suzuki4Exp}, the object
the BCH chain bounds). They are the same operator
(\leanref{Suzuki4Convergence.lean}{371}{suzuki4Step\_eq\_suzuki4Exp}); the four
junction merges of adjacent same-operator exponentials are supplied by the
five-block factorization of \cref{sec:bch-expansion}. Consequently
\eqref{eq:s4-quartic} transfers to \texttt{suzuki4Step}
(\leanref{Suzuki4Convergence.lean}{381}{suzuki4Step\_total\_error\_quartic}:
$\norm{\texttt{suzuki4Step}(n)^n - e^{A+B}} \le C/n^4$, i.e.\ the $t = 1$
instance, since \texttt{suzuki4Step} has the $1/n$ built into its definition;
general $t$ follows by rescaling $A \mapsto tA$, $B \mapsto tB$). The
improvement from $O(1/n^2)$ to $O(1/n^4)$ is therefore recorded on one and the
same object rather than inferred across two definitions.

\subsection{Remaining caveats}
\label{sec:caveats}

Two methodological caveats remain after the Level 4 uniform bound:

\paragraph{Caveat 1: Over-complete basis, non-unique projection.}
The Childs basis of 8 commutators is over-complete (the weight-5 free Lie
algebra is only 6-dimensional). Our projection sets two free parameters to
zero, assigning $\beta_3 = \beta_7 = 0$ (hence $\gamma_3 = \gamma_7 = 0$).
All such projections represent
the \emph{same} $R_5$ element, but the triangle-inequality upper bound
$\sum \abs{\beta_i} \norm{C_i}$ varies. The representation minimizing this
weighted, instance-dependent objective
gives the tightest bound in the Childs basis; finding it requires a linear
program. Our choice is \emph{a} valid projection, not provably \emph{the}
tightest in the Childs basis. It does happen to be strictly smaller than
Childs's on all eight coordinates---the Lean development verifies
$\gamma_i \le \alpha_i$ termwise for the certified ceilings---but that is a property of
\emph{this} projection, not of valid projections in general: adding a large
multiple of a relation vector yields another valid projection with
arbitrarily large coefficients.

\paragraph{Caveat 2: Dependency boundary and generated constants.}
Lean-Trotter carries zero own theorem-level BCH-interface axioms. All
four levels are \emph{fully proved at the project level}:
\texttt{\#print axioms norm\_suzuki4\_childs\_form\_via\_level3},
\texttt{\#print axioms norm\_suzuki4\_level2\_bch}, and \texttt{\#print
axioms norm\_suzuki4\_level3\_bch}, together with
\texttt{\#print axioms norm\_suzuki4\_level4\_uniform}, each return only
\texttt{[propext, Classical.choice, Quot.sound]} at Lean-BCH
revision~\texttt{05e8c52} (2026-07-28). The earlier revision
\texttt{d455ff0} discharged the previous B1.c quintic axiom; the current
pin also discharges the two former septic stepping stones.
The L2, L3 theorems compose, on the Lean-Trotter side, the Lean-BCH
bridge corollaries
\texttt{suzuki5\_log\_product\_quintic\_of\_IsSuzukiCubic} and
\texttt{suzuki5\_log\_product\_quintic\_tight\_at\_suzukiP}, each via
an exp-Lipschitz lift; L1 is a corollary of L3 via the Lean-proved
termwise inequality $\gamma_i \le \alpha_i$. The L4 small-$\tau$
($\tau^7$) refinement composes the bridge corollary
\texttt{suzuki5\_log\_product\_septic\_at\_suzukiP}; its symmetric-BCH
septic bridge and Childs-7-basis matching residual are both theorems in
the pinned companion
\href{https://github.com/Jue-Xu/Lean-BCH}{Lean-BCH} project.
The previously axiomatized
order-4 Taylor-coefficient identity
\texttt{bch\_iteratedDeriv\_s4Func\_order4}, which supports the abstract
$O(t^5)$ existence bound, is now a theorem via a three-slice chain
that composes a single-step $O(\abs{\tau}^5)$ bound from Lean-BCH
(\texttt{norm\_s4Func\_sub\_exp\_le\_of\_IsSuzukiCubic} composed with the
opaque-RHS corollary
\texttt{suzuki5\_bch\_M4b\_RHS\_le\_t5\_of\_IsSuzukiCubic}), a generic
Taylor-match-from-norm lemma
(\leanref{TaylorMatch.lean}{313}{iteratedDeriv\_eq\_of\_norm\_le\_pow}),
and our \texttt{iteratedDeriv\_exp\_smul\_mul\_at\_zero}. Axiom
inspection confirms that this theorem and the other headline results
(Lie--Trotter, Strang, single-step $S_4$ bound, all of L1/L2/L3)
depend only on \texttt{[propext, Classical.choice, Quot.sound]}, with
no transitive \texttt{sorryAx}. The four three-factor
symmetric-BCH-cubic identities previously axiomatized here (the
\texttt{symmetric\_bch\_cubic} family) are likewise theorems derived
from the corresponding Lean-BCH theorems specialized to
$\mathbb{K} = \mathbb{R}$. The $\gamma_i$ and $K$ literals are hard-coded
in the Lean sources and validated there only through the inequalities they
must satisfy (e.g.\ $\abs{\beta_i(\text{suzukiP})} \le \gamma_i$ via
\texttt{nlinarith} on tight rational brackets); the Python computation and
its transcription into Lean are \emph{not} themselves machine-checked.
The Lean theorem checks the inequalities actually used downstream, which is
the relevant trust boundary for the stated bounds.

\subsection{Summary}

This section demonstrates a \emph{CAS-assisted formalization workflow}
producing what is, to our knowledge, the first analytic tightening of the
Childs et al.\ (2021) $S_4$ prefactors (based on a literature search
documented in the project repository). Four bounds are formalized:

\begin{itemize}[leftmargin=*, nosep]
  \item \textbf{L1 (Childs form, axiom-free on Trotter's side).}
    Recovers, locally, the \emph{coefficient form} of Childs et al.'s
    (2021) Prop.~J.1: their coefficients $0.0046$--$0.0284$ with no
    remainder term on some neighbourhood of zero, under the strict-gap
    hypothesis of \cref{tab:levels}. (Their proposition itself is global in
    $t \ge 0$ in the unitary setting; L1 holds in any complete normed
    algebra, but only near $0$.) Derived from L3 via the
    Lean-proved termwise inequality $\gamma_i \le \alpha_i$;
    no axiomatization of Childs et al.'s bound is required.
  \item \textbf{L2 (BCH unit coefficients).} The simplest rigorous
    BCH-derived bound, using unit coefficients on all 8 Childs commutators.
    Follows from the primitive BCH bound $\abs{\beta_i(\text{Suzuki-}p)} \le 1$
    and is dominated numerically by L3.
  \item \textbf{L3 (BCH leading-order).} Replaces Childs et al.'s
    coefficients $\alpha_i$ with certified rational ceilings $\gamma_i$ of
    the signed BCH leading-order coefficients. All 8 are strictly smaller; two ($\gamma_3,
    \gamma_7$) vanish identically, and on the remaining six the improvement
    ratio ranges from $8.6\times$ (for $C_2, C_6$) to $64\times$ (for $C_8$).
    Its sharp form \cref{eq:level3-sharp} carries the $\gamma$-sum in the
    statement; with $\delta, K'$ existential, it certifies the asymptotic
    leading coefficient rather than an end-to-end numerical step count.
  \item \textbf{L4 (BCH small-$\tau$ refinement).} Adds an explicit $R_7$
    correction ($K \approx 0.01951$, also CAS-derived), pushing the first
    correction to the $\tau^5$ term out to order $\tau^7$. Both the
    multiplier $C$ and the radius $\delta$ are existentially quantified, so
    L4 is a demonstration that the CAS pipeline iterates to higher BCH
    order---not a bound one extracts a number from; that role belongs to
    \cref{eq:level3-sharp}.
\end{itemize}

Both the Lean formalization and the CAS computation are reproducible from
the project repository. Levels L1--L4 are fully Lean-verified without
project-specific axioms at Lean-BCH revision \texttt{05e8c52} (see
\cref{sec:caveats}). The three-factor symmetric BCH interface on which
L3/L4 depend is imported from Lean-BCH as a git dependency and is no
longer axiomatized on the Trotter side.
